\section{The ASML Question}\label{sec:background}

\subsection{A chokepoint}

Every chip made at the leading edge, including the accelerators and high-bandwidth memory on which frontier AI runs, is patterned with EUV lithography, and ASML is the only firm that makes EUV machines. It recognized 48 EUV systems in 2025, four of them the newer High-NA version, and expected to ship about 65 low-NA systems in 2026. In October 2026 its chief executive said that planned supply for both 2027 and 2028 would not meet customer demand \citep{asml2026annual, asml2026q3}. The machines cannot simply be bought and left alone. Installed-base management, ASML's term for servicing, upgrades, and spare parts, made up a quarter of its 2025 revenue, and EUV tools reportedly stop working quickly without ASML parts \citep{asml2026annual, bloomberg2024disable}.

China cannot buy EUV machines, but it is ASML's largest market for the older immersion DUV tools with which Chinese fabs make 7-nanometre chips by multiple patterning. Customers in China accounted for 26 percent of ASML's sales in 2023, 36 percent in 2024, and 29 percent in 2025 \citep{asml2026annual}. ASML expects China to fall to about 20 percent in 2026, although Chinese memory makers pulled orders forward in the third quarter, reportedly to stockpile tools against a future cutoff \citep{asml2026q2call, asml2026q3}. Several hundred Dutch immersion tools are believed to be running in China, and they require maintenance roughly every six months \citep{aei2026lithography}. China's own efforts to escape the chokepoint are real but slow. A prototype EUV light source built in Shenzhen with former ASML engineers had not produced working chips by the end of 2025, and Beijing's target of 2028 for working chips is widely considered optimistic \citep{reuters2025euv}.

\subsection{A nested chokepoint}

ASML is a Dutch company, but the machine it sells is not a Dutch product in any simple sense. Carl Zeiss SMT in Germany is the sole supplier of its optics, and ASML says that a prolonged loss of Zeiss supply would ``effectively'' end its business. The EUV light source is developed at Cymer in San Diego and driven by laser amplifiers from Trumpf. At the end of 2025 ASML had 8,242 employees in the United States, about a fifth of its workforce, and its manufacturing sites include Wilton, Connecticut, and San Diego \citep{asml2026annual}. This geography gives Washington jurisdiction over the firm. When the United States restricted support for advanced Chinese fabs by US persons in October 2022, ASML instructed its American staff to stop servicing Chinese customers within a week \citep{cnbc2022asml}.

The United States has turned that jurisdiction into a standing instrument. A December 2024 rule extended the foreign direct product rule to semiconductor manufacturing equipment and set a zero de minimis threshold for certain lithography tools bound for China. It exempted countries with equivalent national controls, among them the Netherlands, Japan, and most EU member states, through a supplement that the United States can amend at will \citep{bis2024fdpr}. In April 2026 a bipartisan bill, the MATCH Act, proposed a countrywide ban on immersion lithography and on servicing for listed Chinese fabs, with 150 days for allies to align before Washington would impose controls of its own \citep{baumgartner2026match}. In the terms of \citet{farrell2019weaponized}, ASML satisfies the network condition for weaponization but not the institutional one: the state that hosts the node does not control the networks in which the node is embedded.

\subsection{Three dependencies}

Europe's use of ASML is shaped by three dependencies, each of which the model treats explicitly.

\paragraph{Security.} Europe relies on the United States for nuclear deterrence, air and missile defence, intelligence, and enablers it cannot quickly replace. Bruegel and the Kiel Institute estimate that defending Europe without the United States would require about 300,000 additional troops and about EUR 250 billion a year in additional spending \citep{burilkov2025defending}. That sum is roughly eight times ASML's 2025 revenue and about 25 times its sales to China. The second Trump administration has made the link between protection and economic alignment explicit. Its December 2025 National Security Strategy promises ``more favorable treatment on commercial matters, technology sharing, and defense procurement'' to partners that ``align their export controls with ours'' \citep{whitehouse2025nss}. It has also cast doubt on the guarantee itself. Withdrawals from Germany were announced in 2026 and a larger drawdown was under review that September, and in a May 2026 poll majorities in all fifteen European countries surveyed doubted the United States would come to their defence \citep{nbc2026pentagon, ecfr2026home}.

\paragraph{Frontier AI.} Europe hosts about five percent of the world's AI compute and depends on American firms for frontier models \citep{juijn2026europe}. Access has already been treated as an instrument of policy. The January 2025 AI Diffusion Rule divided the world into tiers and placed ten EU members in the top tier and the other seventeen below it \citep{bis2025diffusion}. The August 2025 joint statement on the EU--US trade framework tied European purchases of US AI chips to European adoption of US technology-security requirements \citep{whitehouse2025jointstatement}.

\paragraph{China.} Europe imports almost all of its rare-earth permanent magnets from China, and Beijing has shown that it can use licensing to cut supply in weeks. Its October 2025 controls, which reached foreign products containing small amounts of Chinese rare earths, were suspended after the US--China meeting in Busan, not withdrawn \citep{mofcom2025suspension}. The Nexperia dispute is the closest real analogue to the scenario's endgame. On September 30, 2025, the Dutch government used emergency powers to take control of the Chinese-owned chipmaker after a US rule threatened its operations. China banned exports from Nexperia's Chinese plant, European carmakers warned of stoppages within days, and on November 19 the Dutch minister suspended the order \citep{taylorwessing2025nexperia}. Relief came mainly through the US--China deal, not through anything Europe obtained for itself.

\subsection{What the scenario claims}

\emph{Europe 2031} makes five claims that a model can test. (1) In 2028, faced with US secondary sanctions, the Netherlands has no choice but to halt servicing of Chinese DUV tools, and gets nothing for it. (2) In 2029 Europe cannot use ASML against the United States because security-dependent members and the Netherlands itself block the countermeasure. (3) In 2031 Washington moves to take control of ASML because it fears losing the chokepoint to a Europe drifting toward China. (4) Its threat of permanent exclusion from US AI, backed by a stockpile of EUV tools, lets it outlast Europe. (5) Compute on European soil and a coalition of middle powers would have given Europe leverage. The model confirms (1) and (3), rejects the mechanism behind (2), finds the threat in (4) incredible as stated, and specifies when (5) would and would not work.
